\subsection{APPLICATION TO FREE MONOIDS}

Lemma 2. Let \(\mathbf{M}\) be the monoidal sum of the family \((M_x)_{x\in \mathbf{X}}\) defined by \(\mathbf{M}_x = \mathbf{N}\) for all \(x\in \mathbf{X}\) and let \(\phi_{x}\) denote the canonical homomorphism of \(\mathbf{M}_x\) into \(\mathbf{M}\) . The mapping \(x\mapsto \phi_x(1)\) of \(\mathbf{X}\) into \(\mathbf{M}\) extends to an isomorphism \(h\) of \(\operatorname {Mo}(\mathbf{X})\) onto \(\mathbf{M}\) .

Let \(h\) be the homomorphism of \(\operatorname{Mo}(\mathbf{X})\) into \(\mathbf{M}\) characterized by \(h(x) = \phi_x(1)\) . For every integer \(n \geqslant 0\) , \(\phi_x(n) = \phi_x(1)^n = h(x)^n\) and as \(\mathbf{M}\) is generated by \(\bigcup_{x \in \mathbf{X}} \phi_x(\mathbf{N})\) , it is also generated by \(h(\mathbf{X})\) . Hence \(h\) is surjective. Moreover, for all \(x\) in \(\mathbf{X}\) , the mapping \(n \mapsto x^n\) is a homomorphism of \(\mathbf{N} = \mathbf{M}_x\) into \(\operatorname{Mo}(\mathbf{X})\) ; there thus exists (no. 3, Proposition 4) a homomorphism \(h'\) of \(\mathbf{M}\) into \(\operatorname{Mo}(\mathbf{X})\) such that \(h'(\phi_x(n)) = x^n\) for \(x \in \mathbf{X}\) and \(n \in \mathbf{N}\) ; in particular, \(h'(h(x)) = x\) for \(x \in \mathbf{X}\) and hence \(h' \circ h\) is the identity homomorphism of \(\operatorname{Mo}(\mathbf{X})\) . Therefore \(h\) is injective. It has thus been proved that \(h\) is bijective.

PROPOSITION 6. Let \(w\) be an element of \(\mathbf{Mo}(\mathbf{X})\) .

\begin{enumerate}
\def\labelenumi{(\alph{enumi})}
\tightlist
\item
  There exist an integer \(n \geqslant 0\) , elements \(x_{\alpha}\) of \(\mathbf{X}\) and integers \(m(\alpha) > 0\) (for
\end{enumerate}

\(1 \leqslant \alpha \leqslant n\) such that \(x_{\alpha} \neq x_{\alpha + 1}\) for \(1 \leqslant \alpha < n\) and \(w = \prod_{\alpha = 1}^{n} x_{\alpha}^{m(\alpha)}\) . The sequence \((x_{\alpha}, m(\alpha))_{1 \leqslant \alpha \leqslant n}\) is determined uniquely by these conditions.

\[
x _ {\beta} ^ {\prime}
\]

\[
m ^ {\prime} (\beta)
\]

\[
1 \leqslant \beta \leqslant p
\]

(b) \(w = \prod_{\beta=1}^{p} x'^{m'(\beta)}.\) Then \(p \geqslant n\) . If \(p = n\) , then \(x_{\beta}' = x_{\beta}\) and \(m'(\beta) = m(\beta)\) for \(1 \leqslant \beta \leqslant p\) .

In the notation of Lemma 2, \(\bar{h}^{1}(\phi_x(n)) = x^n\) for \(x \in \mathbf{X}\) and \(n \in \mathbf{N}\) . Proposition 6 then follows from no. 3, Proposition 5.

